\section{Rose kernels and dials}

 
 	\subsection{Product-form kernels and minimal assumptions}
	
		%-- Displacement notation ----------------------------------
		\[
   			z := y-x,  \qquad r := \|z\|.
		\]

		\begin{assumptionlist}[label=(W\arabic*)]
			\item \label{W1}(\textbf{Radial weight}) $w:[0,\infty)\to[0,\infty)$ is $C^{1}$, radial, $\int_{\R^{d}} w(\|z\|)\,dz = 1$,  $\int\|z\|^{2} w(\|z\|)\,dz < \infty$, and $w\in L^{2}(\R^{d})$.

			\item \label{W2} (\textbf{Centred residual}) $g:\R^{d}\to[0,\infty)$ is $C^{1}$,  $\int g(z)\,dz = 1$,  $\int z\,g(z)\,dz = 0$, $\int zz^{\top} g(z)\,dz < \infty$,  and $\nabla g\in L^{2}(\R^{d})$.
			
		\end{assumptionlist}
	
		\begin{definition}[Rose kernel]\label{def:RoseKernel}
			For $(\theta,\sigma)\in(0,\infty)^{2}$ set
			\[
				w_\theta(r):=\theta^{-d}w\!\bigl(r/\theta\bigr),  \quad g_\sigma(z):=\sigma^{-d}g\!\bigl(z/\sigma\bigr), \quad k_{\theta,\sigma}(y\mid x):=  \frac{w_\theta(\|x-y\|)\,g_\sigma(y-x)} {\displaystyle\int_X w_\theta(\|x-y'\|)\,g_\sigma(y'-x)\,dy'} 
			\]
		\end{definition}
		
		\begin{lemma}[Normalisation and basic bounds]\label{lem:k-norm-bounds}
			Under Assumptions~\ref{W1}--\ref{W2}
			the expression in Definition~\ref{def:RoseKernel} defines a
			measurable Markov transition density $k_{\theta,\sigma}(y\mid x)$:
			\[
  				0 \le k_{\theta,\sigma}(y\mid x) \le
  				\frac{w_\theta(\|x-y\|)\,g_\sigma(y-x)}{\inf_{x}\!Z_{\theta,\sigma}(x)},
  				\qquad
  				\int_X k_{\theta,\sigma}(y\mid x)\,d\mu(y)=1,
			\]
			where $Z_{\theta,\sigma}(x)$ is the denominator in
			Definition~\ref{def:RoseKernel} and satisfies
			$0<\inf_x Z_{\theta,\sigma}(x)\le \sup_x Z_{\theta,\sigma}(x) < \infty$.
		\end{lemma}

		\begin{assumptionlist}[label=(K\arabic*)]
			\item \label{K0} (\textbf{Markov}) $k_{\theta,\sigma}(y\mid x)\ge0$ and $\int_X k_{\theta,\sigma}(y\mid x)\,dy = 1$.

			\item \label{K1} (\textbf{Score in $L^{2}$}) $S_{\theta,\sigma}(y\mid x):=\nabla_y\log k_{\theta,\sigma}(y\mid x) \in L^{2}(dy\,dx)$.
		
			\item \label{K2} (\textbf{Second-moment scaling}) $\displaystyle\int_X\|y-x\|^{2}k_{\theta,\sigma}(y\mid x)\,dy = O(\theta^{2})$ as $\theta\downarrow 0$.

			\item \label{K3} (\textbf{$C^{1}$ in dials}) $\partial_\theta\log k_{\theta,\sigma},  \partial_\sigma\log k_{\theta,\sigma}\in L^{2}(dy\,dx)$.
		\end{assumptionlist}
	
  	\subsection{Local–linear residual construction (canonical example)}
	
 	\subsection{Scaled kernels and notation for the rest of the paper}
	
	\subsection{Neutral (reversible) Rose operator}
	
		\begin{itemize}[nosep,leftmargin=*]
      			\item Define $k^{\natural}(y\mid x)=\tfrac12[k(y\mid x)+k(x\mid y)]$.
      			\item Self‑adjoint \(\mathcal R^{\natural}\); serves as equilibrium baseline.
    		\end{itemize}

\section{Population theory}
  %-------- assumption block -------------
  	\begin{assumptionbox}
    		(K0)–(K2) as in Sec. 3.0 — Markov, score $L^2$, second-moment scaling.
 	 \end{assumptionbox}

  	\subsection{CGF $\Rightarrow$ Fisher $\Rightarrow$ Landau}
		\begin{definition}[One–step CGF]\label{def:CGF}
			% as in backbone
		\end{definition}
		
		\begin{definition}[Symmetric Fisher load]\label{def:Fisher-load}
			Given a Rose kernel $k_{\theta,\sigma}(y\mid x)$ with score
			$S_{\theta,\sigma}(y\mid x):=\nabla_y\log k_{\theta,\sigma}(y\mid x)$
			and anchor/target measures $\mu_0,\mu_1$,
			define for $f\in W^{1,2}(X,\mu)$
			\[
  				\mathcal I_{\theta,\sigma}[f]
  				\;:=\;
  				\iint_{X\times X}
    				\big\langle \nabla f(y),\,S_{\theta,\sigma}(y\mid x)\big\rangle^{2}
    				\,d\mu_0(x)\,d\mu_1(y).
			\]
			Similarly, for a density $\rho\ll\mu$,
			\[
 				\mathcal I_{\theta,\sigma}[\rho]
  				\;:=\;
 				 \iint_{X\times X}
    				\big\langle \nabla\log\rho(y),\,S_{\theta,\sigma}(y\mid x)\big\rangle^{2}
    				\,d\mu_0(x)\,d\mu_1(y).
			\]
		\end{definition}

		\begin{proposition}[CGF–Fisher correspondence]\label{prop:CGF}
			For every $(\theta,\sigma)\in(0,\infty)^2$ and every centred observable
			$f\in W^{1,2}(X,\mu)$,
			\[
  				\Lambda_{\theta,\sigma}(\lambda;f)
				= \frac{\lambda^{2}}{2}\,\mathcal I_{\theta,\sigma}[f]
    				+ O(\lambda^{3}),
  				\qquad \lambda\to0,
			\]
			where
			\(
  				\mathcal I_{\theta,\sigma}[f]
  		 		= \iint\!\langle\nabla f(y),S_{\theta,\sigma}(y\mid x)\rangle^{2}\,
     				d\mu_0(x)\,d\mu_1(y).
			\)
		\end{proposition}


		\begin{remark}[Legendre dual = Landau stiffness]
			A two–line box reminding readers that the Legendre transform of $\Lambda$ yields a quadratic Landau free energy.
		\end{remark}
	
	\subsection{Dirichlet energy as common source}
     		\begin{definition}[Symmetrised kernel and Dirichlet form]\label{def:Dirichlet}
			Let
				$k^{\natural}_{\theta,\sigma}(y\mid x):=
				\frac12\bigl(k_{\theta,\sigma}(y\mid x)+k_{\theta,\sigma}(x\mid y)\bigr)$
			and let $\mu_\natural$ be any reference measure absolutely continuous
			w.r.t.\ both $\mu_0$ and $\mu_1$ (often $\mu_\natural=\tfrac12(\mu_0+\mu_1)$).
			Define the associated Dirichlet form
			\[
 				 \mathcal E_{\theta,\sigma}(f)
  				:=
  				\frac12
  				\iint_{X\times X}
   			 	\bigl(f(y)-f(x)\bigr)^{2}
    				\,k^{\natural}_{\theta,\sigma}(y\mid x)\,
    				d\mu_\natural(x)\,d\mu_\natural(y).
			\]
		\end{definition}
		
		\begin{lemma}[Fisher load from Dirichlet energy]\label{lem:I-from-E}
			Under Assumptions~\ref{K0}--\ref{K2},
			there exists a dial-dependent coefficient $c_{\theta,\sigma}>0$ such that
			\[
 				\mathcal I_{\theta,\sigma}[f]
  				\;=\;
  				c_{\theta,\sigma}\,\mathcal E_{\theta,\sigma}(f)
 			 	\;+\; o(1)
  				\qquad (\theta\downarrow0),
			\]
			with $c_{\theta,\sigma}
  			= 4/\theta^{2} + o(\theta^{-2})$ in the isotropic small-locality regime.
		\end{lemma}
     
     	\subsection*{Remark: Reversible baseline and operator split}
     	% 5‑line remark splitting $\mathcal R^{\flat}=\mathcal R^{\natural}+A$

  	\subsection{Fisher–curvature–entropy bridge}
  
		\begin{lemma}[Iterated carré–du–champ]\label{lem:curv}
			Under Assumptions~\ref{K0}–\ref{K2}, for every
			$f\in W^{1,2}(X,\mu)$,
			\[
  				\Gamma_{2}^{(\theta,\sigma)}(f)
  				= \frac{\theta^{2}}{2}\,
   				\mathcal I_{\theta,\sigma}[f] + o(\theta^{2}),
  				\qquad \theta\downarrow0 .
			\]
		\end{lemma}

		\begin{lemma}[Discrete de Bruijn Identity]\label{lem:deBruijn}
			Fix $\theta>0$.  
			For every density $\rho\ll\mu$ with $\mathcal I_{\theta,\sigma}[\rho]<\infty$,
			\[
  				\frac{d}{d\sigma}
  				H\!\bigl(R^{\sharp}_{\theta,\sigma}\rho\bigr)
  				= -\frac12\,\mathcal I_{\theta,\sigma}[\rho],
  				\qquad \sigma>0 .
			\]
		\end{lemma}


		\begin{theorem}[Fisher–curvature–entropy bridge]\label{thm:bridge}
			With Assumptions~\ref{K0}–\ref{K2},
			\[
 		 		\boxed{\;
    				\mathcal I_{\theta,\sigma}
    				= \frac{2}{\theta^{2}}\,
        				\Gamma_{2}^{(\theta,\sigma)}
    				= -2\,\frac{d}{d\sigma}
        				H\!\bigl(R^{\sharp}_{\theta,\sigma}\rho\bigr)
 		 		\;}
  				\;+\; o(1).
			\]
		\end{theorem}
		
		\begin{corollary}[Discrete fluctuation--dissipation]\label{cor:FDT}
			Combine Lemma~\ref{lem:curv} and Lemma~\ref{lem:deBruijn}.
			In the small--locality regime $\theta\downarrow0$,
			the response coefficient obtained by differentiating the
			one--step entropy $H(R^{\sharp}_{\theta,\sigma}\rho)$ w.r.t.\ the
			noise dial $\sigma$ equals (up to the universal factor $2/\theta^{2}$)
			the curvature prefactor $\Gamma_{2}^{(\theta,\sigma)}(f)$, i.e.
			\[
  				\frac{d}{d\sigma}H\bigl(R^{\sharp}_{\theta,\sigma}\rho\bigr)
  				\;=\;
  				-\frac{1}{\theta^{2}}\,
   				\Gamma_{2}^{(\theta,\sigma)}(f)
  				\;+\; o(1).
			\]
			Thus fluctuations (Fisher load) determine dissipation (entropy
			production) in exact analogy with the Onsager reciprocal relations.
		\end{corollary}
		
	\subsection{Discrete Onsager--Machlup functional}\label{subsec:OM}
		\begin{corollary}[Discrete OM action]\label{cor:action}
			Let $\mathbf x=(x_0,\dots,x_n)$ be a trajectory sampled every $h$ and
			let $k_{\theta,\sigma}$ satisfy \textnormal{(K0)--(K2)}.
			With local drift $m_\theta$ and quadratic forms
			$\Gamma_t^{2}:=\Gamma^{(\theta,\sigma)}_2(f_t)$,
			$\dot H_t:=\partial_\sigma H(R^{\sharp}_{\theta,\sigma}\rho_t)$,
			\[
  				\mathcal A_{\theta,\sigma}[\mathbf x]
				\;=\;
    				\frac{1}{\theta^{2}}\sum_{t}\Gamma_t^{2}
   				\;-\;
    				\frac{1}{\sigma^{2}}\sum_{t}\dot H_t
    				\;+\; o(1)
			\]
			as $(\theta,\sigma)\!\downarrow0$.  If $k=k^{\natural}$ and
			$\sigma^{2}=2Dh$ with $h\to0$, this reduces to the classical
			continuous Onsager--Machlup functional.
		\end{corollary}

		\begin{remark}[Data‑driven OM]\label{rem:datadrivenOM}
			Any consistent drift estimate (e.g. LWLS–1) may replace $m_\theta$
			without changing the quadratic structure of
			$\mathcal A_{\theta,\sigma}[\mathbf x]$.
		\end{remark}

	%---------------------------------------------------------------------------
	% Deterministic limit & KL ⇔ least action
	%---------------------------------------------------------------------------

	\subsection{Deterministic limit and least–action equivalence}
		\begin{remark}[Classical Lagrangian limit]\label{rem:classical_L}
			In the joint limit $\theta,h\to0$ with $\sigma^{2}=2Dh$, the discrete
			action density converges to
			$\tfrac{1}{4D}\!\int_{0}^{T}\|\dot x(t)-b(x(t))\|^{2}dt$,
			the standard Onsager–Machlup Lagrangian for diffusions.
		\end{remark}

		\begin{theorem}[KL--least–action equivalence]\label{thm:KL_eq_action}
			For every Rose kernel $k_{\theta,\sigma}$ let
				\[
  					J(\theta,\sigma)
  					:= \mathrm{KL}\!\bigl(\mu_1 \,\Vert\, k_{\theta,\sigma}\!\#\mu_0\bigr),
 			 		\qquad
 					 \mathcal A_{\theta,\sigma}^{\mathrm{pop}}
  					:= \E_{\mu_0}\bigl[\mathcal A_{\theta,\sigma}[X_{0:n}]\bigr]
				\]
				denote, respectively, the population Kullback–Leibler functional and the
				population Onsager–Machlup action introduced in
				Corollary~\ref{cor:action}.
				Under Assumptions \textnormal{(K0)--(K3)}
				\[
  					\boxed{\;
    					\nabla_{(\theta,\sigma)} J
    					\;=\;
   					\nabla_{(\theta,\sigma)} \mathcal A_{\theta,\sigma}^{\mathrm{pop}}
  					\;}
 	 				\quad\text{for all }(\theta,\sigma)\in(0,\infty)^2.
				\]
		\end{theorem}
		
		\begin{remark}
		The pair
		\[
  			(\theta^{\star},\sigma^{\star})
  			:=\arg\min_{(\theta,\sigma)} J(\theta,\sigma)
 		 	= \arg\min_{(\theta,\sigma)} \mathcal A_{\theta,\sigma}^{\mathrm{pop}}
		\]
		maximises the likelihood of the observed path ensemble within the two‑dial Rose family; 
		equivalently, it minimises the expected Onsager–Machlup action. Hence trajectories 
		generated under the little‑rose kernel are the most probable (in the maximum‑likelihood 
		sense) among all trajectories that can be produced by Rose kernels.
		\end{remark}


 	 \subsection{Introducing the dial plane}
		\begin{assumptionbox}
 			 \textbf{(K4) C$^{1}$ regularity in dials.}
			$(\theta,\sigma)\mapsto\log k_{\theta,\sigma}$ is $C^{1}$ and
			$\partial_\theta\log k_{\theta,\sigma},
			\partial_\sigma\log k_{\theta,\sigma}\in L^{2}(d\mu_0d\mu_1)$.
		\end{assumptionbox}
	
	\subsection{Edge‑wise dial oracle}  % ← moved up
		\begin{assumptionbox}[frametitle={Oracle axioms}]
			\begin{enumerate}[label=\textbf{(A\arabic*)},leftmargin=*]
 		 		\item \emph{Reversibility consistency.}
       					If $k(y\mid x)=k(x\mid y)$ then ${\cal O}(x_i,x_j)={\cal O}(x_j,x_i)$.
 	 			\item \emph{Permutation equivariance.}
      					${\cal O}(\pi(i),\pi(j))={\cal O}(x_i,x_j)$ for any permutation $\pi$.
  				\item \emph{Strict‑convex minimal risk.}
        					${\cal O}(x_i,x_j)=\arg\min_{(\theta,\sigma)}\mathcal J_{ij}(\theta,\sigma)$ with $\mathcal J$ strictly convex.
				\item\label{A4}\textbf{Positive empirical mass.}  
      					Empirical anchor weights satisfy $\displaystyle\inf_{i\in V}\,\widehat{\mu}_0(i)\;>\;0$.
			\end{enumerate}
		\end{assumptionbox}
		
		\begin{definition}[Edge KL functional]\label{def:edge-KL}
			Given observed transitions from anchor $x_i$,
			let $\widehat{k}_{ij}$ denote the empirical conditional probability that
			a step from $x_i$ lands in $x_j$.
			Define the edge-wise KL at dials $(\theta,\sigma)$:
			\[
  				\mathcal J_{ij}(\theta,\sigma)
  				:=
  				\sum_{j}
    				\widehat{k}_{ij}\,
    				\log\frac{\widehat{k}_{ij}}{k_{\theta,\sigma}(x_j\mid x_i)}.
			\]
			(For continuous $X$, replace the sum by an integral against the
			empirical conditional measure $\widehat{k}_{i}$.)
		\end{definition}
		
		\begin{proposition}[Edge‑wise KL gradient]\label{prop:edgeKLgrad}
			For every ordered pair $(x_i,x_j)$ and all $(\theta,\sigma)$,
			\[
  				\nabla_{(\theta,\sigma)}
  				\mathcal J_{ij}(\theta,\sigma)
  				= -\frac12\,
    				\nabla_{(\theta,\sigma)}
    				\mathcal I_{ij}(\theta,\sigma).
			\]
			Consequently the oracle dials
			$(\theta^{\star}_{ij},\sigma^{\star}_{ij})$ maximise the Fisher load
			and minimise the edge‑wise KL.
		\end{proposition}


		\begin{definition}[Little‑rose operator]\label{def:littleRose}
			Assign each edge its oracle dials $(\theta_{ij}^{\star},\sigma_{ij}^{\star})$ and denote by $\mathcal R_{\star}^{\flat/\sharp}$ the resulting pull‑back and push‑forward operators.  We call this pair the \emph{little‑rose} operator.
		\end{definition}

	\subsection{Global KL gradient corollary}
		\begin{corollary}[Global KL gradient]\label{cor:globalKL}
			Summing Proposition~\ref{prop:edgeKLgrad} over all edges with weights $\mu_0(i)$ yields
			\[
  				\nabla_{(\theta,\sigma)}
  				\mathrm{KL}\!\bigl(\mu_1
        				\parallel k_{\theta,\sigma}\!\#\mu_0\bigr)
 				= -\;\tfrac12\,
   				\nabla_{(\theta,\sigma)}
    				\mathcal I_{\theta,\sigma},
			\]
			hence the little‑rose operator is the unique minimiser of KL and maximiser of the Fisher load.
		\end{corollary}

\section{Finite-sample dial selection}
 	\subsection{Goldenshluger–Lepski selector}
	
		\begin{lemma}[Bias--variance decomposition]\label{lem:bv}
			For either PF or Koopman risk,
			\[
 	 			\E\,\widehat{\mathcal R}^{\square}(\theta,\sigma)
 			 	\;=\;
  				\mathrm{Bias}^{2}_{\square}(\theta,\sigma)
  				\;+\;
  				\mathrm{Var}_{\square}(\theta,\sigma),
			\]
			where the bias term is the squared $L^{2}$ distance between the
			population Rose operator at $(\theta,\sigma)$ and the “oracle” operator
			that best fits the data class, and the variance term scales as
			\[
  				\mathrm{Var}_{\square}(\theta,\sigma)
  				\;\lesssim\;
  				\frac{ \mathrm{trace}\bigl(\Sigma_{\theta,\sigma}^{\square}\bigr) }{n_{\mathrm{eff}}(\theta)},
			\]
			with an effective sample size $n_{\mathrm{eff}}(\theta)$ determined by the
			local weight $w_\theta$ (see Appendix~\ref{app:GL}).
		\end{lemma}
  
  		\begin{theorem}[GL selector: non‑asymptotic risk]\label{thm:GL}
			Let $\widehat{\theta},\widehat{\sigma}$ be chosen by the GL rule
			with penalty parameter $\alpha$. Then with probability at least
			$1-\delta$,
			\[
  				\mathcal R\bigl(\widehat{\theta},\widehat{\sigma}\bigr)
 				\;\le\;
  				C_{\alpha}\,\inf_{(\theta,\sigma)\in\mathcal G}
      				\mathcal R(\theta,\sigma)
  				\;+\;
  				C\frac{\log(1/\delta)}{n},
			\]
			where $\mathcal R$ denotes either PF or Koopman $L^{2}$‑risk and
			$C_{\alpha},C$ are explicit constants (see Appendix~\ref{app:GL}).
		\end{theorem}
		
		\begin{definition}[Sample risks]\label{def:sample-risk}
			Let $\widehat{R}^{\sharp}_{\theta,\sigma}$ and
			$\widehat{R}^{\flat}_{\theta,\sigma}$ denote the empirical PF and
			Koopman estimators built from the sample.
			Given a test class $\mathcal F\subset L^{2}(X,\mu)$ and
			$\mathcal P\subset L^{1}(X,\mu)$ define
			\[
				\widehat{\mathcal R}^{\mathrm{PF}}(\theta,\sigma)
  				:= \sup_{\rho\in\mathcal P}
       				\big\| \widehat{R}^{\sharp}_{\theta,\sigma}\rho
            			- R^{\sharp}_{\theta,\sigma}\rho \big\|_{L^{2}}^{2},
 				 \qquad
  				\widehat{\mathcal R}^{\mathrm{K}}(\theta,\sigma)
  				:= \sup_{f\in\mathcal F}
       				\big\| \widehat{R}^{\flat}_{\theta,\sigma}f
            			- R^{\flat}_{\theta,\sigma}f \big\|_{L^{2}}^{2}.
			\]
		\end{definition}

\section{Directed curvature and reversibility diagnostics}
	\subsection{Curvature proxy on the dial graph}
	\subsection{Irreversibility measure via antisymmetric current}
	
		\begin{definition}[Antisymmetric weight \& current]\label{def:Aij}
			For a Rose kernel $k_{\theta,\sigma}$ define the log--ratio antisymmetry
			\[
  				A_{ij} := \frac12
 		 		\Bigl(
    					\log k_{\theta,\sigma}(x_j\mid x_i)
    					- \log k_{\theta,\sigma}(x_i\mid x_j)
 				\Bigr)
			\]
			Given a reference mass $\pi_i>0$ on vertices (e.g.\ stationary or
			empirical mass), the probability current is
			\[
  				J_{ij} := \pi_i k_{\theta,\sigma}(x_j\mid x_i)
         	 		- \pi_j k_{\theta,\sigma}(x_i\mid x_j),
			\]
			satisfying $J_{ij}=-J_{ji}$.
			Detailed balance holds iff all $A_{ij}=0$ (equivalently $J_{ij}=0$).
		\end{definition}
		
		\begin{lemma}[Graph Helmholtz split]\label{lem:graph-Helmholtz}
			Let $J$ be an antisymmetric edge flow on a finite connected graph
			$(V,E)$ with vertex weights $\pi$.
			Then there exists a potential $\phi:V\to\R$ and a cycle flow
			$h\in \ker B^{\top}$ (where $B$ is the weighted incidence matrix)
			such that
			\[
  				J_{ij} = \pi_i k^{\natural}_{ij}(\phi_j-\phi_i) + h_{ij},
			\]
			with $\sum_{j} \pi_i k^{\natural}_{ij}(\phi_j-\phi_i)=0$ at stationarity.
			The decomposition is unique up to an additive constant in $\phi$.
			Moreover $h=0$ iff detailed balance holds.
		\end{lemma}
	
		\begin{remark}[Loschmidt symmetry and the Fisher clock]
 			Setting every antisymmetric weight to zero, $A_{ij}=0$, collapses the
 			kernel to its neutral form $k^{\natural}$ and makes the Rose operator
  			exactly time–reversal symmetric, matching Loschmidt’s objection to
  			irreversible entropy growth.

  			Turning on any $A\neq0$ injects a directed cycle curvature
  			$F=\sum_{\text{cycle}}A$, and the Fisher–curvature–entropy bridge
  			(Thm.~\ref{thm:bridge}) yields
 			 \[
     				F \;=\; \frac{1}{2}\,\mathcal I_{\theta,\sigma}\;=\; \text{entropy production per step}.
  			\]
  			A non‑zero curvature therefore acts as an intrinsic \emph{Fisher
  			clock}—time emerges only once detailed balance is broken.  In the
  			Loschmidt‑symmetric limit $A\!\to\!0$ the clock halts and the macroscopic
  			arrow of time disappears.
		\end{remark}